\documentclass[10pt,a4paper]{amsart}
\usepackage[margin=19mm]{geometry}
\usepackage{amsmath,amssymb,mathtools}
\usepackage[scr=rsfs,cal=euler]{mathalfa}
\usepackage{xfrac}
\usepackage[unicode,hidelinks,bookmarks=false]{hyperref}
\newcommand{\ie}{i.e.\ }
\newcommand{\cf}{cf.\ }
\newcommand{\resp}{respectively\ }
\newcommand{\Subsection}{\subsection}
\providecommand{\nobreakdash}{\nobreak\hbox{-}}
\setlength{\emergencystretch}{2em}
\multlinegap0pt
\usepackage{tikz}
\usetikzlibrary{cd}
\usepackage{amssymb}
\usepackage{dsfont}
\usepackage{bm}
\usepackage{enumerate}
\newtheorem{proposition}{Proposition}
\newtheorem{theorem}{Theorem}
\newcommand{\Q}{\mathbb{Q}}
\newcommand{\Z}{\mathbb{Z}}
\DeclareMathOperator{\ch}{char}
\newcommand{\rng}[1]{\langle #1 \rangle}
\title{When does an infinite ring have a finite compressed commuting graph?}
\author{Ivan-Vanja Boroja, Damjana Kokol Bukov\v{s}ek, Nik Stopar}
\date{Source: arXiv:2411.07358v1 (CC BY 4.0)}
\begin{document}
\maketitle

\noindent\emph{Source and licence.} When does an infinite ring have a finite compressed commuting graph?. Version arXiv:2411.07358v1 (CC BY 4.0), distributed by arXiv under the Creative Commons Attribution 4.0 International licence, http://creativecommons.org/licenses/by/4.0/, as recorded in the arXiv licence metadata for this exact version and shown on the abstract page https://arxiv.org/abs/2411.07358v1. Full licence notice: \texttt{licenses/arxiv-2411.07358-CC-BY-4.0.txt}.\par\smallskip

\noindent\textit{Editorial scope.} Counted unit: the article's proposition thm:inf\_ring\_1 (source0.tex lines 453--455). The article's complete original proof of it is delivered with it (source0.tex lines 457--531, one \texttt{proof} environment of 75 lines). No mathematics is written, repaired or re-derived: the statement and the proof are the article's own lines, in the article's own notation, and no retained line is substituted or re-flowed.  Retained uncounted as the source-local context the proof needs: lines 213--215; lines 278--280.  Nothing was omitted from any retained span, so the package carries no omission record.  No retained line contains a citation, so the wrapper carries no bibliography and no bibliography entry.  The retained lines contain the article's own diagram code, which the wrapper typesets with the standard tikz packages; no image file is delivered and the package declares no asset dependency.  Editorial formatting only, in addition: an AMS \texttt{amsart} preamble that restates the article's own declarations and macros used by the retained lines, namely the environments \texttt{proposition}, \texttt{theorem} on separate counters, together with the standard packages those lines need.  The retained environments print in this document as Proposition 1, Theorem 1 in order of appearance, whereas the article prints the same items with the numbering of its own full document.  The complete native source of the article as distributed in the arXiv e-print for this exact version is bundled unchanged as \texttt{sources/168366/source0.tex}.
\begin{proposition}\label{prop:inf_ring_1}
    If $R$ is an infinite unital ring, then $|V(\Lambda^1(R))|=|R|$, unless $|R|=\aleph_0$,  $\ch R = 0$, and every element of $R$ is annihilated by some nonzero linear polynomial with coefficients in $\Z$.
\end{proposition}

\begin{theorem}\label{thm:inf_ring}
    If $R$ is an infinite ring, then $|V(\Lambda(R))|=|R|$.
\end{theorem}

\begin{proposition}\label{thm:inf_ring_1}
    If $R$ is an infinite unital ring, then either $|V(\Lambda^1(R))|=|R|$ or $R$ is isomorphic to a unital semidirect product $\Z[\frac1m] \ltimes I$ for some positive integer $m$ and some finite $\Z[\frac1m]$-ring $I$.
\end{proposition}

\begin{proof}
Assume $R$ is an infinite unital ring such that $|V(\Lambda^1(R))| \neq |R|$.
Proposition~\ref{prop:inf_ring_1} implies that $|R|=\aleph_0$ (hence, $|V(\Lambda^1(R))|<\infty$), $\ch R = 0$, and every element of $R$ is annihilated by some nonzero linear polynomial with coefficients in $\Z$.

The above allows us to define a map $g \colon R \to \Q$ by $g(r)=\frac{a}{b}$ if $r$ is annihilated by the linear polynomial $bx-a \in \Z[x]$.
The fact that $\ch R = 0$ implies $b \neq 0$ for all $r \in R$.
The image $g(r)$ is independent of the choice of the linear polynomial that annihilates $r$. Indeed, if $b_1 x-a_1$ and $b_2x-a_2$ both annihilate $r$, then $b_2a_1=b_1b_2r=b_1a_2$, so $\frac{a_1}{b_1}=\frac{a_2}{b_2}$.
In addition, $g$ is a unital ring morphism, namely, if $r,s \in R$, $bx-a$ annihilates $r$ and $dx-c$ annihilates $s$, then $bd(r+s)=da+bc$ and $bd(rs)=ac$, so $g(r+s)=g(r)+g(s)$ and $g(rs)=g(r)g(s)$. Clearly, $g(1)=1$.
Hence, $g(R)$ is a unital subring of $\Q$.

Clearly, $|V(\Lambda^1(g(R)))| \leq |V(\Lambda^1(R))|<\infty$, which implies that $g(R)$ is finitely generated as a unital ring.
Hence, $g(R)=\Z[\frac{a_1}{b_1},\frac{a_2}{b_2},\ldots,\frac{a_n}{b_n}]$ for some $a_i,b_i \in \Z$, where $a_i$ and $b_i \neq 0$ are relatively prime for each $i$.
It follows that $g(R)=\Z[\frac1m]$, where $m$ is the least common multiple of $b_1,b_2,\ldots,b_n$.
So we have a surjective unital ring morphism $g \colon R \to \Z[\frac1m]$.

Let $I$ be the kernel of $g$ and denote the canonical inclusion of $I$ into $R$ by $f$. Then we have an exact sequence
\begin{equation}\label{eq:exact}
\begin{tikzcd}[sep=20pt]
0 & I & R & \Z[\frac1m] & 0 \ .
\arrow[from=1-1, to=1-2]
\arrow["f", from=1-2, to=1-3]
\arrow["g", from=1-3, to=1-4]
\arrow[from=1-4, to=1-5]
\end{tikzcd}
\end{equation}
Since $\ch R = 0$, we may regard $\Z$ as a subring of $R$.
So $\Z+I$ is a unital subring of $R$. Since $g$ is unital, its restriction to $\Z$ is the identity map so that $\Z \cap I=0$. 

If $k \in \Z$ and $a \in I$, the ring $\rng{k+f(a)}_1$ is generated by $k+f(a)$ and $1$, thus also by $f(a)$ and $1$.
Hence, it contains $\Z\cdot 1 +\rng{f(a)}=\Z + \rng{f(a)} \subseteq \Z+I$. But the latter is clearly a unital subring of $R$, therefore $\rng{k+f(a)}_1=\Z + \rng{f(a)}=\Z\cdot 1 + f(\rng{a})=\rng{f(a)}_1$.
This implies that the map $\hat{f} \colon \Lambda(I) \to \Lambda^1(\Z+I)$ given by $\hat{f}([a])=[f(a)]_1$ is a well defined bijection on the sets of vertices.
Furthermore, for any $a,b \in I$ we have $ab=ba$ if and only if $f(a)f(b)=f(b)f(a)$. Hence, $\hat{f}$ is a graph isomorphism.
This implies that $\Lambda^1(\Z+I) \cong \Lambda(I)$.
Hence, $|V(\Lambda(I))|=|V(\Lambda^1(\Z+I))| \leq |V(\Lambda^1(R))|<\infty$ and the ring $I$ is finite by Theorem~\ref{thm:inf_ring}.

Choose an element $r \in R$ such that $g(r)=\frac1m$. Then $mr-1 \in I$.
Since $(mr-1)^1$, $(mr-1)^2$, $(mr-1)^3,\ldots \in I$ and $I$ is finite, we infer $(mr-1)^u=(mr-1)^v$ for some positive integers $u>v$.
This implies that $mr$ is integral of degree $\leq u$.
In addition, $mr-1$ must have finite additive order, say $n$, so that $n(mr)=n \in \Z$.
Putting the two observations together we infer
$\rng{mr} \subseteq \Z+M$ where
$$M=\Big\{n_1(mr)+n_2(mr)^2+\ldots+n_{u-1}(mr)^{u-1} ~|~ 0 \leq n_1,n_2,\ldots,n_{u-1} <n\Big\}.$$
Since $(mr)^2,(mr)^{2^2},(mr)^{2^3},\ldots \in \rng{mr}$ and the set $M$ is finite we must have $(mr)^{2^j}-(mr)^{2^i} \in \Z$ for some positive integers $j>i$.
On the other hand $(mr)^{2^j}-(mr)^{2^i} \in I$, since $g((mr)^{2^j})=1=g((mr)^{2^i})$, so we conclude that $(mr)^{2^j}-(mr)^{2^i}=0$, because $\Z \cap I=0$.
This implies that the element $e_0=(mr)^k$ with $k=2^j-2^i>0$ is an idempotent in $R$. Indeed,
$$e_0^2=(mr)^{2k}=(mr)^{2^j+(2^j-2\cdot 2^i)}=(mr)^{2^j}(mr)^{2^j-2\cdot 2^i}=(mr)^{2^i}(mr)^{2^j-2\cdot 2^i}=(mr)^{2^j-2^i}=e_0,$$
since $2^j-2\cdot 2^i \geq 0$ (here and in the rest of the proof it should be understood that in a calculation a factor with exponent $0$ is omitted).

Now let $s=r(mr)^{2k-1}$ and define a subring $S=\rng{s}$. We have $g(s)=\frac1m$. Note that $e_0=e_0^2=ms$, hence $e_0 \in S$.
In addition, $s=r(mr)^{2k-1}=(r(mr)^{k-1})e_0=e_0(r(mr)^{k-1})$, so $se_0=s=e_0s$, thus $e_0$ is an identity element in $S$.

We want to prove that $S \cap I=0$. Suppose $a \in S \cap I$.
Then $a=p(s)$ for some polynomial $p \in \Z[x]$ with $p(0)=0$.
Furthermore, applying $g$ to $a$ we obtain $p(\frac1m)=0$.
Denote $d=\deg p$ and define a polynomial $P(x)=x^dp(\frac1x) \in \Z[x]$.
Then $P(m)=0$, so $P(x)=H(x)(x-m)$ for some $H \in \Z[x]$ with $\deg H < \deg P \leq \deg p=d$.
We thus obtain $p(x)=x^dP(\frac1x)=x^{d-1}H(\frac1x)(1-mx)=h(x)(1-mx)$, where $h(x)=x^{d-1}H(\frac1x) \in \Z[x]$.
Furthermore, $h(0)=0$ since $p(0)=0$, so $h(x)=q(x)x$ for some $q \in \Z[x]$ and $p(x)=q(x)(1-mx)x$.
Thus, we get $a=p(s)=q(s)(1-ms)s=q(s)(1-e_0)s=q(s)(s-s)=0$. So $S \cap I=0$.

This implies that $g|_S$, the restriction of $g$ to $S$, is injective. Furthermore, it is also surjective since $g(s)=\frac1m$ and $\Z[\frac1m]$ is generated by $\frac1m$ as a ring, and $g|_S$ is unital since $g(e_0)=g(ms)=mg(s)=1$.
So $g|_S \colon S \to \Z[\frac1m]$ is a unital ring isomorphism.
Define $h \colon \Z[\frac1m] \to R$ by $h(w)=(g|_S)^{-1}(w)$. Then $h$ is a unital ring morphism that extends the exact sequence \eqref{eq:exact} to a commuting diagram
$$\begin{tikzcd}[sep=20pt]
& & & \Z[\frac1m] & \\
0 & I & R & \Z[\frac1m] & 0 \ .
\arrow[from=2-1, to=2-2]
\arrow["f", from=2-2, to=2-3]
\arrow["g", from=2-3, to=2-4]
\arrow[from=2-4, to=2-5]
\arrow["h"', from=1-4, to=2-3]
\arrow["id", from=1-4, to=2-4]
\end{tikzcd}$$
We conclude that $R \cong \Z[\frac1m] \ltimes I$.
\end{proof}

\end{document}
